% L12a student companion deck: synthetic markets with jump hidden Markov models.
% Markov chains and hidden Markov models (from the 2025 lecture text), the jump hidden
% Markov model of one asset, and the sector factor model of many stocks.
%
% November 3, 2026. The frame order mirrors the lecture notebook. The disclaimer is
% slide 2, by instructor preference. Notation follows the lecture: state space
% \mathcal{S}, transition matrix \mathbf{P} with entries p_{ij}, estimate \hat{p}_{ij},
% stationary distribution \bar{\pi}, emission matrix \mathbf{E}, daily growth rate g_t
% (1/yr), jump probability \epsilon and mean episode length \lambda, and the sector
% factor model g_i = alpha_i + beta_i g_m + delta_i f_s + s_i eps_i.
\documentclass[aspectratio=169,11pt]{beamer}
\usepackage{vnslides}

\newcommand{\Sc}{\mathcal{S}}
\newcommand{\Pb}{\mathbf{P}}
\newcommand{\repo}{https://github.com/varnerlab/CHEME-5660-CourseRepository-Fall-2026/blob/main/lectures/week-12/L12a}

\title{\texorpdfstring{{\vnheavy L12a}}{L12a}: Synthetic Markets with Jump Hidden Markov Models}
\subtitle{CHEME 5660 Cornell University}
\author{Copyright Jeffrey Varner 2026. All Rights Reserved}

\begin{document}

\vntitlepage

\begin{frame}{Disclaimer and Risks}
  \vspace{0.6em}
  \textbf{\ink{Training only.}} This content is for training and informational
  purposes only. The teaching team does not make, give, or endorse any offer or
  solicitation to buy or sell securities or derivative products, or provide
  investment or trading advice or strategy.

  \vspace{1.1em}
  \textbf{\ink{Risk.}} Trading involves risk and is generally not recommended for
  those with limited resources, experience, or a low risk tolerance. Past
  performance does not guarantee future success.

  \vspace{1.1em}
  \textbf{\ink{Responsibility.}} You are responsible for your investment decisions
  based on your financial circumstances, objectives, risk tolerance, and
  liquidity needs.
\end{frame}

\begin{frame}{Objectives for Today}
  Today, we build synthetic markets: simulated price histories with the features of
  real ones. How would a trading method do on thousands of plausible paths?

  \vspace{0.5em}
  {\large\textbf{\ink{Today's concepts}}}
  \begin{itemize}\setlength{\itemsep}{0.4em}
    \item \hilite{Markov chains and hidden Markov models}: states, transition
      matrices, stationary distributions, emissions.
    \item \hilite{A jump hidden Markov model of one asset}: binned states, counted
      transitions, Student-t emissions, jump episodes.
    \item \hilite{A synthetic market of many stocks}: shared market and sector draws
      plus each stock's own draw.
  \end{itemize}
\end{frame}

\begin{frame}{What You Need to Know}
  \vspace{0.3em}
  {\large\textbf{\ink{Core concepts}}}
  \begin{itemize}\setlength{\itemsep}{0.4em}
    \item A Markov chain is a set of states $\Sc$ and a transition matrix $\Pb$ whose
      rows sum to one; its stationary distribution satisfies $\bar{\pi}=\bar{\pi}\Pb$.
    \item The jump hidden Markov model bins the daily growth rate into states,
      estimates $\Pb$ by counting, and draws each day from a Student-t of the day's
      state. Jump episodes (probability $\epsilon$, mean length $\lambda$) keep
      volatility clustering.
    \item A synthetic stock is
      $g_{i}=\alpha_{i}+\beta_{i}\,g_{m}+\delta_{i}\,f_{s}+s_{i}\,\tilde{\varepsilon}_{i}$:
      shared market and sector draws make stocks move together.
  \end{itemize}

  \vspace{0.5em}
  {\large\textbf{\ink{For understanding}} \dimmed{(you won't reproduce these)}}
  \begin{itemize}\setlength{\itemsep}{0.3em}
    \item Power iteration, the tuning objective for $\epsilon$ and $\lambda$, and how
      the sector loadings and own-draw scale are chosen.
  \end{itemize}
\end{frame}

\begin{frame}{Lectures and Examples}
  \vspace{0.3em}
  {\normalsize\textbf{\ink{Lecture:}}\par\small
  \href{\repo/CHEME-5660-L12a-Lecture-JumpHMM-SyntheticMarkets-Fall-2026.ipynb}{\dimmed{CHEME-5660-L12a-Lecture-JumpHMM-SyntheticMarkets-Fall-2026.ipynb}}}

  \vspace{0.6em}
  {\normalsize\textbf{\ink{Examples:}}\par\small
  \href{\repo/CHEME-5660-L12a-Example-JumpHMM-SPY-Fall-2026.ipynb}{\dimmed{CHEME-5660-L12a-Example-JumpHMM-SPY-Fall-2026.ipynb}}\par
  \href{\repo/CHEME-5660-L12a-Example-JumpHMM-SectorFactor-Fall-2026.ipynb}{\dimmed{CHEME-5660-L12a-Example-JumpHMM-SectorFactor-Fall-2026.ipynb}}}
\end{frame}

% ========== Why Synthetic Markets?
\begin{frame}{Why Synthetic Markets?}
  A backtest runs a method on one historical path, and a method that did well may
  have been lucky. A \hilite{synthetic market} generates many plausible paths. For
  stress tests, its paths need:
  \begin{itemize}\setlength{\itemsep}{0.4em}
    \item \hilite{Fat tails}: large daily moves are far more common than a normal
      distribution predicts.
    \item \hilite{Volatility clustering}: large moves follow large moves, so $|g_{t}|$
      stays correlated for weeks.
    \item \hilite{Co-movement}: stocks move together, through the market and their
      sectors.
  \end{itemize}
  \vspace{0.3em}
  Geometric Brownian motion has none of the three.
\end{frame}

% ========== Markov Chains
\begin{frame}{Markov Chains}
  A discrete Markov model $\mathcal{M}=(\Sc,\Pb)$ has a state space $\Sc$, the set of
  values the system can take, and a transition matrix $\Pb$. It has the
  \hilite{Markov property}: given the present state, the past does not matter:
  \begin{equation*}
    P(X_{n+1}=s_{n+1}\mid X_{0}=s_{0},\dotsc,X_{n}=s_{n}) = P(X_{n+1}=s_{n+1}\mid X_{n}=s_{n})
  \end{equation*}
  The one-step transition probability from $s_{i}$ to $s_{j}$ is
  $p_{ij}=P(X_{n+1}=s_{j}\mid X_{n}=s_{i})$, the same for every $n$ (time-homogeneous).
\end{frame}

\begin{frame}{The Transition Matrix}
  The $p_{ij}$ fill the transition matrix $\Pb\in\mathbb{R}^{|\Sc|\times|\Sc|}$: rows
  are current states, columns next states, $p_{ij}\geq 0$, and each row sums to one.
  For weather with $\Sc=\{\text{Sunny},\text{Rainy}\}$:
  \begin{equation*}
    \Pb=\begin{bmatrix}0.8 & 0.2\\ 0.4 & 0.6\end{bmatrix}
    \qquad
    \pi_{1}=\pi_{0}\Pb=[1,\,0]\begin{bmatrix}0.8 & 0.2\\ 0.4 & 0.6\end{bmatrix}=[0.8,\,0.2]
  \end{equation*}
  If today is Sunny, tomorrow is Sunny with probability 0.8. The distribution
  $\pi$ is a row vector: $\pi^{\prime}_{j}=\sum_{i}\pi_{i}p_{ij}$.
\end{frame}

\begin{frame}{Estimating the Transition Matrix}
  We do not know $\Pb$; we estimate it from observed states $X_{1},\dotsc,X_{T}$. With
  $n_{ij}$ the number of days in $s_{i}$ followed by a day in $s_{j}$:
  \begin{equation*}
    \hat{p}_{ij}=\frac{n_{ij}}{\sum_{k=1}^{|\Sc|}n_{ik}}
  \end{equation*}
  the share of days in $s_{i}$ that moved to $s_{j}$. Each row sums to one by
  construction; this is the maximum likelihood estimate, and how a market
  generator estimates its transition matrix.
\end{frame}

\begin{frame}{The Stationary Distribution}
  The \hilite{stationary distribution} $\bar{\pi}$ is unchanged by a step,
  $\bar{\pi}=\bar{\pi}\Pb$: the long-run share of time in each state, and where a
  generator starts its paths. \hilite{Power iteration} applies $\Pb$ until nothing changes:
  \begin{itemize}\setlength{\itemsep}{0.25em}
    \item Start from the uniform $\pi=[1/N,\dotsc,1/N]$.
    \item Repeat $\pi^{\prime}\gets\pi\Pb$ until $\lVert\pi^{\prime}-\pi\rVert_{1}\leq\text{tol}$ or
      $T$ iterations; return $\bar{\pi}\gets\pi^{\prime}$.
  \end{itemize}
  \vspace{0.2em}
  This needs an ergodic $\Pb$ (irreducible, aperiodic); the spectral gap sets the speed.
\end{frame}

% ========== Hidden Markov Models
\begin{frame}{Hidden Markov Models}
  In a \hilite{hidden Markov model} $\mathcal{M}=(\Sc,\mathcal{O},\Pb,\mathbf{E})$ we do
  not see the states, only what they emit: observations $\mathcal{O}$ with emission
  probabilities $e_{ij}=P(O_{n}=o_{j}\mid X_{n}=s_{i})$, rows summing to one. To sample:
  draw $X_{0}\sim\bar{\pi}$; then each step draws $X_{n}$ from row $X_{n-1}$ of $\Pb$
  and $O_{n}$ from row $X_{n}$ of $\mathbf{E}$.

  \vspace{0.5em}
  \textbf{\ink{Continuous emissions:}} a growth rate is a number, so each state has an
  emission \emph{distribution} instead of a row of $\mathbf{E}$:
  $g_{n}\mid X_{n}=s_{k}\sim$ emission distribution of $s_{k}$.
\end{frame}

% ========== The Jump Hidden Markov Model
\begin{frame}{The Jump Hidden Markov Model}
  For the daily growth rate $g_{t}=\frac{1}{\Delta t}\ln(S_{t}/S_{t-1})$ (1/yr) of one asset:
  \begin{itemize}\setlength{\itemsep}{0.35em}
    \item \hilite{States}: fit a Laplace distribution to the growth rates and cut it into
      $N$ bins of equal probability; state 1 holds the most negative growth rates.
    \item \hilite{Transitions}: assign each day to its bin and count, $\hat{p}_{ij}$.
    \item \hilite{Emissions}: in state $s_{k}$,
      $g_{t}\mid X_{t}=s_{k}\sim\mu_{k}+\sigma_{k}\,T_{\nu}$, a Student-t located and scaled
      by the days in that state.
  \end{itemize}
  \vspace{0.2em}
  In the data the states are observed bins; in simulation they are hidden, and only
  the growth rates are reported.
\end{frame}

\begin{frame}{Why Add Jumps?}
  The counted chain forgets quickly: after a volatile day, a few transitions usually
  return it to the middle states, so $|g_{t}|$ is correlated for only a few days.
  \hilite{Jump episodes} keep volatile periods going:
  \begin{itemize}\setlength{\itemsep}{0.35em}
    \item With no episode under way, an episode starts with probability $\epsilon$.
    \item It lasts a Poisson number of days with mean $\lambda$.
    \item Each episode day is drawn from the few most extreme states, negative or
      positive: a run of large moves of both signs.
  \end{itemize}
\end{frame}

\begin{frame}{Tuning and Checking the Model}
  A grid search picks $(\epsilon,\lambda)$. On simulated paths that contain an
  episode, it scores the autocorrelation $\rho_{k}$ of $|g_{t}|$ and the kurtosis
  $\kappa$ against the data:
  \begin{equation*}
    J(\epsilon,\lambda)=\sum_{k=1}^{25}\left(\rho_{k}-\hat{\rho}_{k}\right)^{2}+0.2\left(\kappa-\hat{\kappa}\right)^{2}
  \end{equation*}
  where hats mark a simulated path, averaged over the paths. The tuned model is
  checked on the distribution, the kurtosis, and the volatility clustering.
  \slidenote{Example:}{\href{\repo/CHEME-5660-L12a-Example-JumpHMM-SPY-Fall-2026.ipynb}{\dimmed{CHEME-5660-L12a-Example-JumpHMM-SPY-Fall-2026.ipynb}}}
\end{frame}

% ========== Many Assets
\begin{frame}{Many Assets: The Sector Factor Model}
  Each synthetic stock $i$ in sector $s$ is built from draws it shares with other
  stocks and a draw of its own:
  \begin{equation*}
    g_{i}(t)=\alpha_{i}+\beta_{i}\,\underbrace{g_{m}(t)}_{\textstyle\text{market}}+\delta_{i}\,\underbrace{f_{s}(t)}_{\textstyle\text{sector}}+s_{i}\,\underbrace{\tilde{\varepsilon}_{i}(t)}_{\textstyle\text{own}}
  \end{equation*}
  $\alpha_{i}$, $\beta_{i}$: single index estimates, held fixed. $f_{s}$: a model of the
  sector's average residual, with loading $\delta_{i}$. $\tilde{\varepsilon}_{i}$: a model of
  the stock, centered and scaled by $s_{i}$ so the stock keeps its real variance. Each
  piece is a jump hidden Markov model.
\end{frame}

\begin{frame}{What the Model Keeps and Gives Up}
  Composing stocks from shared draws on the same day has consequences:
  \begin{itemize}\setlength{\itemsep}{0.45em}
    \item \hilite{Keeps}: each stock's fat tails and volatility clustering; stocks in a
      sector move together through their shared sector draw, and all stocks through
      the market draw.
    \item \hilite{Gives up}: sector draws are independent, so stocks in different
      sectors are linked only through the market, and cross-sector correlations come
      out a little low.
    \item \hilite{No copula}: a copula reorders each stock's draws across days and would
      scramble the timing of its volatile periods and jump episodes.
  \end{itemize}
\end{frame}

\begin{frame}{An Option: A Drifting $\beta$}
  Holding $\beta_{i}$ fixed is a simplification: real $\beta$ wanders. The model can let
  $\beta_{i}(t)$ drift around $\hat{\beta}_{i}$ as an exponential moving average of random
  shocks, sized by how much the online estimate of $\beta$ wanders in the real data.

  \vspace{0.5em}
  A trading method that re-estimates $\beta$ as it trades then faces the same problem
  in the synthetic market as in the real one.
  \slidenote{Example:}{\href{\repo/CHEME-5660-L12a-Example-JumpHMM-SectorFactor-Fall-2026.ipynb}{\dimmed{CHEME-5660-L12a-Example-JumpHMM-SectorFactor-Fall-2026.ipynb}}}
\end{frame}

\begin{frame}{Summary}
  \vspace{0.3em}
  We built synthetic markets for stress tests, from Markov chains to a jump hidden
  Markov model of one asset and a sector factor model of many stocks.

  \vspace{0.4em}
  {\large\textbf{\ink{Key takeaways}}}
  \begin{itemize}\setlength{\itemsep}{0.4em}
    \item \hilite{Markov chains}: states, a transition matrix estimated by counting,
      and a stationary distribution; hidden states emit what we observe.
    \item \hilite{Jump hidden Markov model}: fat tails from binned Student-t states, and
      part of the volatility clustering from tuned jump episodes.
    \item \hilite{Sector factor model}: shared market and sector draws make synthetic
      stocks move together.
  \end{itemize}
  \vspace{0.3em}
  Synthetic markets let us test a method on many paths, not one.
\end{frame}

\end{document}
